\section{结论}

本文将结构反事实对比学习（\scc）作为 GRIP Stage~2 的关系判别扩展：保留参数内化流程、原始提示和生成式 QA 目标，新增关系候选级 InfoNCE，并以训练图关系画像和冻结 GRIP 基线分数构建可追溯候选库。单种子主实验报告了验证选定配置相对本地 generation-only 基线的直接解码结果，但配置搜索限制了对新增目标贡献的单独归因。固定配置对照进一步表明，结果随候选组成、数据集和答案规则变化：FB 中 Full 与频率配额匹配组接近，NELL 的直接解码与重排差异方向相反，WN18RR 的历史条件对照则呈现较大正差异，但受实体对核验缺口和基线复现差异限制。现有证据支持在 GRIP 中评估关系候选级对比目标，并通过频率分层、候选数量和分布诊断刻画其表现边界；它尚未确立普遍收益或独有采样机制。联合报告生成表现、候选排序、候选来源及其频率组成，可以更准确地界定这项 GRIP 扩展。
